\section{Arhitectura Sistemelor Distribuite și Backend}

Pentru a integra componenta de inteligență artificială într-o aplicație accesibilă utilizatorilor finali, StyleGenAI adoptă o arhitectură de sistem distribuit, bazată pe modelul Client-Server. Această secțiune analizează principiile arhitecturale care au guvernat proiectarea sistemului backend.

\subsection{Principiile Arhitecturale REST}
Interfața de programare a aplicației (API) a fost dezvoltată respectând constrângerile stilului arhitectural \textbf{REST (Representational State Transfer)}, definit de Roy Fielding. Aceste constrângeri asigură scalabilitatea și independența componentelor:

\begin{itemize}
    \item \textbf{Arhitectura Client-Server}: Există o separare clară între interfața utilizator (Client - aplicația mobilă) și logica de stocare/procesare (Server - backend-ul Python). Aceasta permite evoluția independentă a celor două componente.
    \item \textbf{Stateless (Fără Stare)}: Fiecare cerere HTTP de la client către server trebuie să conțină toate informațiile necesare pentru a fi înțeleasă și procesată. Serverul nu păstrează starea sesiunii utilizatorului între cereri, ceea ce simplifică designul și permite scalarea orizontală (adăugarea de noi instanțe de server fără a sincroniza sesiuni).
    \item \textbf{Interfață Uniformă}: Resursele (Utilizatori, Haine, Ținute) sunt identificate prin URI-uri (Uniform Resource Identifiers) și manipulate prin metode HTTP standard (GET, POST, PUT, DELETE).
\end{itemize}

\subsection{Ecosistemul Python pentru Calcul Științific}
Deși limbaje precum Node.js sau Go sunt populare pentru microservicii web datorită modelului lor de concurență asincronă, Python a fost ales pentru acest proiect datorită supremației sale în domeniul calculului numeric și al inteligenței artificiale.
Framework-ul \textbf{Flask} servește drept "adeziv" între cererile HTTP și bibliotecile de procesare intensivă:
\begin{itemize}
    \item \textbf{Interoperabilitate}: Flask permite apelarea directă a funcțiilor din \texttt{PyTorch} și \texttt{scikit-learn} în cadrul aceluiași proces, eliminând latența de comunicare inter-procesuală care ar apărea dacă AI-ul ar fi un serviciu separat.
    \item \textbf{WSGI (Web Server Gateway Interface)}: Aplicația respectă standardul PEP 3333, permițând rularea sa pe orice server de aplicații compatibil (ex: Gunicorn, uWSGI) în mediul de producție.
\end{itemize}

\subsection{Design-ul Modular al Serviciilor}
Intern, aplicația backend este structurată pe module funcționale, respectând principiul \textit{Separation of Concerns}:
\begin{enumerate}
    \item \textbf{Stratul de Rutare (Controller)}: Gestionează parsing-ul cererilor HTTP și validarea datelor de intrare.
    \item \textbf{Stratul de Servicii (Business Logic)}: Implementează algoritmii specifici (ex: logica de generare a ținutelor, orchestrarea procesării imaginilor).
    \item \textbf{Stratul de Acces la Date (DAO/Repository)}: Abstractizează interacțiunea cu baza de date, permițând schimbarea tehnologiei de stocare fără a afecta restul aplicației.
\end{enumerate}
Această arhitectură stratificată facilitează testarea unitară și mentenanța pe termen lung a codului.